% SPDX-License-Identifier: Apache-2.0
% covers: Wdog
\datasheet{Wdog}{A watchdog timer: a countdown software has to keep refreshing, and a reset when it stops.}

\dsfacts{\code{lib/parts/src/wdog.rs}}{\code{txhdl\_parts::wdog::Wdog}}%
{\code{//docs:examples}}

\subsection*{Function}

A program that hangs stays hung until somebody presses the button, and
on a board nobody is standing next to, nobody does. So the hardware
counts down and software has to keep saying that it is still there.
When it stops saying so, \code{Wdog} asks the system control block of
\code{Syscon} for a reset, which records the watchdog as the cause, so
a program that restarts can read why it did.

Three things make it a watchdog rather than a timer, and each is there
because a hung program can do more than stop.

A refresh carries \code{KEY}. A program that has lost its way is
writing whatever it finds over whatever it reaches, and a watchdog it
can refresh by accident is a watchdog it cannot fail. A wrong key is
not ignored: it is a failure.

In window mode a refresh while the count is still above \code{sill} is
too early, and too early is also a failure. A program looping tightly
on its own refresh is as stuck as one that has stopped, and a plain
countdown cannot tell them apart.

The lock sets and does not clear. After it is set, \code{ctrl},
\code{load} and \code{sill} are read only until the next reset.

\subsection*{Parameters}

\begin{center}\footnotesize
\begin{tabular}{@{}lp{0.68\textwidth}@{}}
\toprule
Parameter & Meaning \\
\midrule
\code{KEY} & The word a refresh carries. A write of anything else to
\code{feed} is a failure rather than a refresh. \\
\bottomrule
\end{tabular}
\end{center}

\subsection*{Register map}

Offsets from the block's base. The block selects the word by bits 4 to
2 of the address and looks at no other bits.

\dsregs{Wdog}

Once \code{lock} is set, writes to \code{ctrl}, \code{load} and
\code{sill} are ignored, and the lock bit itself only ever sets. A write
of \code{KEY} to \code{feed} refreshes the count, and a write of
anything else is a failure; \code{feed} reads zero. In window mode, a
refresh while \code{count} is above \code{sill} is a failure too.

The field is \code{sill} and not \code{open}, because \code{open} is a
reserved word of VHDL and the lowering refused such a name when the
part was written (issue 77). It now escapes one instead, to
\code{open\_rw} (issue 497), and the field keeps the name it has.

\dstables{Wdog}

\subsection*{Behaviour}

\begin{itemize}
\item While \code{ctrl} bit 0 is set, \code{count} falls by one a
cycle. A refresh loads it from \code{load}.
\item A write of \code{KEY} to \code{feed} is a refresh, unless window
mode is on and \code{count} is above \code{sill}, in which case it is
a failure. A write of anything else is a failure in either mode.
\item When \code{count} reaches zero: if warning is on and has not
happened, the warning bit sets, \code{irq} rises and the count
reloads, which gives a program that is merely slow one more timeout to
notice. Otherwise it is a failure.
\item A failure sets the failed bit, drives \code{rst\_req} and holds it
for eight cycles, and reloads the count. The reload matters: the count
is at zero when a timeout fails, so a watchdog that did not reload
would ask for a reset on every cycle from then on and hold the machine
down rather than restart it.
\item \code{irq} is a level, standing while the warning bit is set and
warning is on. \code{rst\_req} is a request to \code{Syscon}, not a reset
line.
\item Nothing in this block is reset by its own request, which is what
lets a program read \code{status} after the restart.
\end{itemize}

\subsection*{Verification}

\begin{itemize}
\item \code{//lib/examples:ex\_wdog} sets a timeout, watches the count
fall between two reads and come back after a refresh, stops refreshing
and takes the warning and then the reset, writes a wrong key, refreshes
while the window is shut, and locks the watchdog on and fails to turn
it off.
\item The build simulates the netlist against that run under nvc and
Verilator: \code{//docs:wdog\_sim\_wdog\_tb\_test} and
\code{//docs:wdog\_vsim\_test}.
\end{itemize}

\subsection*{Used by}

Nothing in the tree yet. \code{Syscon} has the \code{wdog} line this
drives, and issue 149 asks for a board program that stops refreshing
and reads \code{watchdog} as the cause after it restarts.

\subsection*{Limits}

\code{count} and \code{load} are sixteen bits, so the timeout is at
most 65535 cycles. The reset request is eight cycles and not a
parameter. \code{feed} reads as zero rather than as what was last
written. There is no second stage: the warning and the reset are the
only two things a timeout can do.
